\documentclass[11pt,a4paper]{article}
\usepackage[margin=17mm]{geometry}
\usepackage[T1]{fontenc}
\usepackage{lmodern,amsmath,amssymb,xcolor,tabularx,microtype}
\usepackage[hidelinks]{hyperref}
\definecolor{Ink}{HTML}{16324F}
\pagestyle{empty}
\setlength{\parindent}{0pt}
\setlength{\parskip}{3pt}
\setlength{\emergencystretch}{1em}
\newcommand{\question}[2]{\vspace{7pt}{\color{Ink}\bfseries #1. #2}\par}
\newcommand{\choice}[2]{\par\textbf{#1.}\ #2\par}
\hypersetup{pdftitle={Quiz: Lecture 7},pdfauthor={Scientific Computing 2026}}
\begin{document}
{\sffamily\color{Ink}\LARGE Quiz: Lecture 7}\par
{\sffamily Scientific Computing 2026 \hfill 10 minutes \quad 5 questions}
\vspace{4pt}\hrule\vspace{6pt}
Name: \rule{9cm}{0.4pt}\hfill Date: \rule{2.5cm}{0.4pt}

Choose one answer per question.

\question{1}{Entropy}
Consider two distributions:
\[
P=(0.5,0.5),\qquad Q=(0.9,0.1).
\]
How do their entropies compare?

\choice{A}{$H(P)=H(Q)$, because both distributions have two possible outcomes.}
\choice{B}{$H(P)>H(Q)$, because the outcome under $P$ is less predictable.}
\choice{C}{$H(P)<H(Q)$, because one outcome under $Q$ is rare.}
\choice{D}{$H(P)=0$, because its probabilities are equal.}

\question{2}{Cross-entropy and KL divergence}
Suppose the entropy of $P$ is $H(P)=0.6$ and the cross-entropy is
$H(P,Q)=0.9$, using natural logarithms. What is $D_{\mathrm{KL}}(P\|Q)$?

\begin{tabularx}{\linewidth}{@{}l@{\qquad}l@{\qquad}l@{\qquad}X@{}}
\textbf{A.}\ $-0.3$ &
\textbf{B.}\ $1.5$ &
\textbf{C.}\ $0.3$ &
\textbf{D.}\ It cannot be calculated without knowing $H(Q)$.
\end{tabularx}

\question{3}{Distillation}
A teacher provides a fixed distribution $P$. A student predicts $Q_\theta$,
where $\theta$ denotes its parameters. Consider two training objectives:
\[
L_1(\theta)=H(P,Q_\theta),\qquad
L_2(\theta)=D_{\mathrm{KL}}(P\|Q_\theta).
\]
Which statement is correct?

\choice{A}{The objectives differ by a constant independent of $\theta$, so they have the same minimizing parameter values.}
\choice{B}{The objectives have exactly the same value for every $\theta$.}
\choice{C}{They have the same minimizers only if the student can reproduce $P$ exactly.}
\choice{D}{Minimizing either objective encourages the student to assign probability one to the teacher's most likely outcome.}

\question{4}{At the likelihood maximum}
For an observed dataset, the log-likelihood is
\[
\ell(\theta)=C-8(\theta-2)^2,
\]
where $C$ does not depend on $\theta$. At the maximum likelihood estimate,
what are the score and observed information?

\begin{tabularx}{\linewidth}{@{}XX@{}}
\textbf{A.}\ Score $=0$, information $=0$. &
\textbf{B.}\ Score $=16$, information $=0$.\\[3pt]
\textbf{C.}\ Score $=0$, information $=8$. &
\textbf{D.}\ Score $=0$, information $=16$.
\end{tabularx}

\question{5}{Efficient estimation}
Independent Gaussian observations have unknown mean $\mu$ and known variance.
The total Fisher information about $\mu$ is $25$. An unbiased estimator of
$\mu$ has variance $0.04$. What can we conclude?

\choice{A}{Such an estimator cannot exist: its variance must be at least $1/\sqrt{25}=0.2$.}
\choice{B}{It attains the Cram\'er--Rao bound and is efficient among unbiased estimators.}
\choice{C}{Every biased estimator must also have variance at least $0.04$.}
\choice{D}{The estimator equals the true mean for every possible sample.}
\end{document}
